\documentclass[11pt,a4paper]{article}

\usepackage[T1]{fontenc}
\usepackage[utf8]{inputenc}
\usepackage[spanish,es-noshorthands]{babel}
\usepackage{lmodern}
\usepackage{amsmath,amssymb}
\usepackage[margin=2.2cm]{geometry}
\usepackage{microtype}
\usepackage{enumitem}
\usepackage[colorlinks=true,linkcolor=black,citecolor=black,urlcolor=blue]{hyperref}

\newcommand{\dd}{\,\mathrm{d}}
\newcommand{\p}{\partial}
\newcommand{\F}{\mathcal{F}}
\newcommand{\MER}{\mathrm{MER}}
\newcommand{\Rv}{\mathcal{R}_v}

\title{\vspace{-1.3cm}
\textbf{Propuesta: problema inverso en un conducto bifásico 2D}\\[0.4em]
\large Conservación de masa y momentum, $\varphi=\varphi(r,z)$}
\author{Ignacio Gómez Gómez\\[0.15em]
{\small Doble titulación Geología / Ingeniería Matemática}}
\date{\today}

\begin{document}
\maketitle
\vspace{-0.8cm}

\noindent
Propuesta de trabajo. El directo 1D bifásico ya corre. Se pasa a
\textbf{2D estacionario} (sin tiempo), axisimétrico: $\varphi$
también cambia con el radio. Estacionario \textbf{no} apaga la
inercia: solo quita $\p/\p t$. El inverso, para empezar, es un
escalar en la base.

% ---------------------------------------------------------------------------
\section*{1. Dominio y variables}

Cilindro $\Omega=[0,R]\times[H,0]$ ($z$ hacia arriba, $H<0$ es la
base). Flujo estacionario, axisimétrico, sin swirl. En general
\[
P=P(r,z),\qquad
\varphi=\varphi(r,z),\qquad
\mathbf{u}_m=(u_{m,r},u_m),\qquad
\mathbf{u}_g=(u_{g,r},u_g).
\]
$u_m$ y $u_g$ son las componentes verticales; $u_{m,r}$ y $u_{g,r}$
las radiales. Fundido: $\rho_m$ constante. Gas: $\rho_g=P/(\Rv T)$.
En la base se quieren recuperar la sobrepresión $\Delta P$ (entra
en $P|_{z=H}$) y, si cabe después, el agua total $c_0$.

% ---------------------------------------------------------------------------
\section*{2. Conservación de masa y $\varphi=\varphi(r,z)$}

$\Gamma$ es la tasa de exsolución (fundido $\to$ gas).
\begin{align}
\label{eq:masa-m}
\frac{1}{r}\frac{\p}{\p r}\bigl(r\,\rho_m(1-\varphi)\,u_{m,r}\bigr)
+\frac{\p}{\p z}\bigl(\rho_m(1-\varphi)\,u_m\bigr)
&=-\Gamma,
\\[0.25em]
\label{eq:masa-g}
\frac{1}{r}\frac{\p}{\p r}\bigl(r\,\rho_g\varphi\,u_{g,r}\bigr)
+\frac{\p}{\p z}\bigl(\rho_g\varphi\,u_g\bigr)
&=\Gamma.
\end{align}
El caudal másico total no depende de $z$:
\begin{equation}
\label{eq:Q}
Q
=\int_0^R 2\pi r\Bigl[
\rho_m(1-\varphi)\,u_m
+\rho_g\,\varphi\,u_g
\Bigr]\,\dd r
=\mathrm{const}.
\end{equation}
Eso es el $\MER$. Henry:
\begin{equation}
\label{eq:henry}
n(P;c_0)
=\max\Bigl(0,\;
\frac{(1-\xi_0)c_0-C_1 P^{\beta}}{1-C_1 P^{\beta}}\Bigr).
\end{equation}
En cada $(r,z)$ la fracción de gas del flujo local es esa $n$,
\begin{equation}
\label{eq:n-local}
n
=\frac{\rho_g\,\varphi\,u_g}{\rho_m(1-\varphi)\,u_m+\rho_g\,\varphi\,u_g},
\end{equation}
y se despeja
\begin{equation}
\label{eq:phi-r}
\varphi(r,z)
=\frac{n(P;c_0)\,\rho_m\,u_m}
{n(P;c_0)\,\rho_m\,u_m
+(1-n(P;c_0))\,\rho_g\,u_g}.
\end{equation}
Como $u_m$ se anula en la pared y $P$ (luego $n$) puede cambiar
con $r$, $\varphi$ no es plano en el radio.

% ---------------------------------------------------------------------------
\section*{3. Momentum: inercia, $P(r,z)$ y $F_{mw}$}

Estacionario $\Rightarrow$ no hay $\p\mathbf{u}/\p t$. El término
convectivo $(\mathbf{u}\cdot\nabla)\mathbf{u}$ \textbf{sí} está.
En el 1D que ya corre (Kozono--Koyaguchi) es
$\rho_m(1-\varphi)\,u_m\,\dd u_m/\dd z$ y
$\rho_g\varphi\,u_g\,\dd u_g/\dd z$. En 2D, eje $z$:
\begin{align}
\label{eq:mom-m-z}
\rho_m(1-\varphi)\Biggl(
u_{m,r}\frac{\p u_m}{\p r}+u_m\frac{\p u_m}{\p z}
\Biggr)
&=
-(1-\varphi)\frac{\p P}{\p z}
-\rho_m(1-\varphi)\,g
+F_{mg}
\nonumber\\
&\quad
+\frac{1}{r}\frac{\p}{\p r}\Biggl[
r\,\mu(\varphi,P)\,(1-\varphi)\frac{\p u_m}{\p r}
\Biggr]
+\frac{\p}{\p z}\Biggl[
\mu(\varphi,P)\,(1-\varphi)\frac{\p u_m}{\p z}
\Biggr],
\\[0.45em]
\label{eq:mom-g-z}
\rho_g\varphi\Biggl(
u_{g,r}\frac{\p u_g}{\p r}+u_g\frac{\p u_g}{\p z}
\Biggr)
&=
-\varphi\frac{\p P}{\p z}
-\rho_g\varphi\,g
-F_{mg}
\nonumber\\
&\quad
+\frac{1}{r}\frac{\p}{\p r}\Biggl[
r\,\mu_g\,\varphi\frac{\p u_g}{\p r}
\Biggr]
+\frac{\p}{\p z}\Biggl[
\mu_g\,\varphi\frac{\p u_g}{\p z}
\Biggr].
\end{align}
$F_{mg}=F_{mg}(\varphi,u_g-u_m,r_b)$ es el arrastre entre fases
(cuerpo; sobre el gas va con el signo contrario). La fricción de
pared \textbf{no} es un cuerpo: es el borde del viscoseo, abajo.

\paragraph{La presión sí puede depender de $r$.}
Eso lo decide el momentum radial (fundido; el gas es análogo):
\begin{equation}
\label{eq:mom-m-r}
\begin{aligned}
\rho_m(1-\varphi)\Biggl(
u_{m,r}\frac{\p u_{m,r}}{\p r}+u_m\frac{\p u_{m,r}}{\p z}
\Biggr)
&=
-(1-\varphi)\frac{\p P}{\p r}
-\mu(1-\varphi)\frac{u_{m,r}}{r^2}
\\
&\quad
+\frac{1}{r}\frac{\p}{\p r}\Biggl[
r\,\mu(1-\varphi)\frac{\p u_{m,r}}{\p r}
\Biggr]
+\frac{\p}{\p z}\Biggl[
\mu(1-\varphi)\frac{\p u_{m,r}}{\p z}
\Biggr].
\end{aligned}
\end{equation}
Si $u_{m,r}\not\equiv 0$ o el viscoseo radial no es despreciable,
$\p P/\p r\neq 0$ y $P=P(r,z)$. No se impone $P=P(z)$ de entrada.

\paragraph{Reducción de lubricación (primer ataque numérico).}
Si $\varepsilon=R/|H|\ll 1$ (Calbuco: $16/7000$),
\eqref{eq:mom-m-r} da $\p P/\p r=O(\varepsilon)$ y se toma
$P=P(z)$. En el fundido el Reynolds es chico: la inercia de
\eqref{eq:mom-m-z} es $\ll$ el viscoseo radial, y
$\p/\p z$ del esfuerzo axial también. Queda el Stokes radial que
ya se discretiza en el 2D actual. En el gas la inercia vertical
se \textbf{conserva} (como en el 1D). Eso es una reducción, no
el sistema: el de mates es~\eqref{eq:masa-m}--\eqref{eq:mom-m-r}.

\paragraph{$F_{mw}$ y $F_{gw}$ en derivadas.}
Integrando el viscoseo radial del fundido en la sección, el
término de borde es la fricción parietal por unidad de volumen:
\begin{equation}
\label{eq:Fmw}
F_{mw}
=-\frac{2}{R}
\Biggl[
\mu(\varphi,P)\,(1-\varphi)
\frac{\p u_m}{\p r}
\Biggr]_{r=R}.
\end{equation}
($\p u_m/\p r|_{r=R}<0$, así que $F_{mw}>0$: la pared frena.)
En el 1D se reemplaza por $c_g\mu u_m/R^2$; aquí no. Análogo
para el gas cuando toca la pared (tramo fragmentado):
\begin{equation}
\label{eq:Fgw}
F_{gw}
=-\frac{2}{R}
\Biggl[
\mu_g\,\varphi
\frac{\p u_g}{\p r}
\Biggr]_{r=R}.
\end{equation}
Antes de fragmentar, $F_{gw}=0$ (el gas no es la fase continua
en la pared). En el PDE 2D, $F_{mw}$ y $F_{gw}$ \textbf{no} se
meten como cuerpos: el no-deslizamiento ya está en el borde, y
\eqref{eq:Fmw}--\eqref{eq:Fgw} son la identidad al promediar
en $r$ (para comparar con el 1D).

% ---------------------------------------------------------------------------
\section*{4. Condiciones de borde}

Todas las que se conocen. $\p\Omega$ son eje, pared, base y boca.

\paragraph{Eje $r=0$ (simetría / regularidad).}
\begin{equation}
\label{eq:bc-eje}
u_{m,r}=u_{g,r}=0,
\qquad
\frac{\p u_m}{\p r}=\frac{\p u_g}{\p r}=0,
\qquad
\frac{\p \varphi}{\p r}=0,
\qquad
\frac{\p P}{\p r}=0.
\end{equation}

\paragraph{Pared $r=R$ (impermeable, no deslizamiento del fundido).}
\begin{equation}
\label{eq:bc-pared}
u_m=0,
\qquad
u_{m,r}=0,
\qquad
u_{g,r}=0.
\end{equation}
El fundido no resbala. El gas: o $u_g=0$ (no deslizamiento; entonces
\eqref{eq:Fgw} es la tensión), o una ley de deslizamiento si se
quiere más tarde. Sin flujo de masa a través de la pared: las
componentes radiales de~\eqref{eq:masa-m}--\eqref{eq:masa-g} se
anulan en $r=R$. $\varphi$ en la pared no se impone: sale
de~\eqref{eq:phi-r} con $u_m=0$.

\paragraph{Base $z=H$ (cámara).}
\begin{equation}
\label{eq:bc-base}
P(r,H)=P_{\mathrm{lit}}+\Delta P
=\rho_{\mathrm{crust}}g|H|+\Delta P,
\qquad
\xi=\xi_0.
\end{equation}
Si $P(H)$ está bajo saturación, $\varphi(r,H)=0$ y
$u_m(r,H)=u_g(r,H)$ (un solo fluido). Si ya hay gas,
$\varphi(r,H)$ sale de~\eqref{eq:phi-r} con $n(P(r,H);c_0)$.
El perfil $u_m(\,\cdot\,,H)$ tiene que cumplir~\eqref{eq:bc-eje}
y $u_m(R,H)=0$; su promedio fija $Q$ (o al revés: se tira con
$v_{\mathrm{in}}$ y $Q=\pi R^2\rho_{t,i}v_{\mathrm{in}}$).
$v_{\mathrm{in}}$ \textbf{no} es dato: es el parámetro de tiro
del directo, igual que en 1D.

\paragraph{Boca $z=0$ (salida).}
El sistema es el mismo; cambia el cierre de salida (como en 1D).
Efusivo: llega sin fragmentar.
Explosivo: fragmentó, y en la boca se cumple \textbf{una} de
dos cosas (Yoshida--Koyaguchi): o la presión ya bajó a la
atmosférica, o el gas llegó a su velocidad del sonido
isoterma $c_s=\sqrt{\Rv T}$ (flujo sónico: no puede seguir
acelerando dentro del conducto; $P(0)$ puede quedar
$>P_{\mathrm{atm}}$). No es una onda de choque.
\begin{align}
\label{eq:bc-boca}
\text{efusivo:}&\quad
P(r,0)=P_{\mathrm{atm}}
\quad\text{y}\quad
\varphi\le\varphi_{\mathrm{crit}},
\\
\text{explosivo:}&\quad
P(r,0)=P_{\mathrm{atm}}
\quad\text{o}\quad
u_g(r,0)\approx c_s=\sqrt{\Rv T}.
\end{align}
Eso cierra el tiro: se busca $v_{\mathrm{in}}$ (luego $Q$) para
cumplir~\eqref{eq:bc-boca}.

% ---------------------------------------------------------------------------
\section*{5. El problema inverso (para empezar)}

Dado el dato en la base, se resuelve
\eqref{eq:masa-m}--\eqref{eq:mom-m-r} con~\eqref{eq:bc-eje}--\eqref{eq:bc-boca}
y se lee $Q$ en la boca.
\[
\F(\theta)=Q,
\qquad
\theta=\Delta P
\quad\text{o}\quad
\theta=c_0.
\]
Dato: $Q^{\mathrm{obs}}$. Recuperar primero $\Delta P$ (es la
condición de borde~\eqref{eq:bc-base}). Después, si cabe, $c_0$
en~\eqref{eq:henry}. No los dos a la vez.

En mates: inverso de una condición de borde en un sistema 2D
estacionario de conservación. Un escalar. En 2D, $u_m(r)$ y
$\varphi(r,z)$ cambian $Q$; el $\Delta P$ recuperado no es el
del 1D.

No se recupera $R(z)$, no hay transiente, no se reescribe el
solver 1D.

% ---------------------------------------------------------------------------
\section*{6. Qué se le pide al profesor}

¿Tiene sentido como tesis de DIM este inverso (un parámetro, EDP
estacionaria 2D de masa + momentum, con inercia, $P=P(r,z)$ y
$\varphi=\varphi(r,z)$)? Si la lubricación $P=P(z)$ le parece el
primer paso correcto, se parte de ahí. Primero $\Delta P$,
después $c_0$.

\end{document}
